\section{Consequences and remaining distinctions}\label{sec:consequences}
\subsection{Updated energy and negative edits}
The choice of reference energy changes the threshold law. For SPD
$B=A+uu^T$, set
\[
 D=B^{-\alpha/2}(B^\alpha-A^\alpha)B^{-\alpha/2}.
\]
If $q_j$ are the relative eigenvalues of $(B^\alpha,A^\alpha)$,
then the eigenvalues of $D$ are $1-q_j^{-1}$. Indeed the generalized
eigenvalues of $(A^\alpha,B^\alpha)$ are the reciprocals of those
of $(B^\alpha,A^\alpha)$. Hence, for $0<\varepsilon<1$, the exact
unrestricted strict-count supremum at fixed leverage $\tau>0$ is
\[
 \max\left\{0,\left\lceil
 \frac{\alpha\log(1+\tau)}{-\log(1-\varepsilon)}
 \right\rceil-1\right\}.
\]
The weighted-path construction preserves sharpness after changing the
reference energy in this way.

This also describes a conductance decrease that preserves connectivity.
If $L_-=L-dbb^T$ and $\eta=db^TL^\dagger b<1$, reversing the pair
gives a positive addition of leverage $\tau=\eta/(1-\eta)$.
The original relative correction is negative semidefinite, and the
strict count of its singular values above $0<\varepsilon<1$ has
unrestricted weighted-path supremum
\[
 \max\left\{0,\left\lceil
 \frac{-\alpha\log(1-\eta)}{-\log(1-\varepsilon)}
 \right\rceil-1\right\}.
\]
Reversing the unweighted insertion examples gives a deletion with
$\eta=1/2$. These structural statements follow by reversing endpoints;
Theorem~\ref{alg:thm:main} is stated for positive edits. Removing a bridge
changes the kernel and lies outside the common-support setting.

\subsection{Budget-preserving spectral convergence}
The unweighted transfer proves more than a bound on the largest extra
eigenvalue. Use the notation of Lemma~\ref{u:lem:transfer}, and put
$k=d-1$, $T=\alpha\log(1+\tau)$, and
\[
 \delta_s=T-\sum_{j=1}^k\log\mu_{s,j}\longrightarrow0.
\]
The full relative eigenvalues are $\lambda_{s,j}$; pad the macro list
$q_j$ by ones. Every logarithm lies in $[0,T]$. The budget identity and
Ritz lower bounds give
\[
 \sum_{j=1}^k(\log\lambda_{s,j}-\log\mu_{s,j})
       +\sum_{j>k}\log\lambda_{s,j}=\delta_s.
\]
Since the exponential function is $e^T$-Lipschitz on this interval,
\[
 \sum_j|\lambda_{s,j}-q_j^{\rm padded}|
 \le\sum_{j=1}^k|\mu_{s,j}-q_j|+e^T\delta_s\longrightarrow0.
\]
Thus the full-budget graph sequence can be chosen to converge in
$\ell^1$ for correction eigenvalues. For the fixed-dimensional stages
before this transfer, ordinary continuity gives the same convergence;
sequential choices combine them. The weighted unit-mass path transfer
also has this property: for each fixed block size its high responses
vanish as the internal conductance grows, and that conductance can be
chosen after the finite enlarged dimension is known.

The finite-target uniform closure with a strict budget deficit is a
different assertion. If the target has logarithmic budget $T-S$,
$S>0$, append $m$ entries $e^{S/m}$. The maximum added correction tends
to zero, but the added logarithmic sum remains $S$. Such sequences do
not converge to that finite target in $\ell^1$. This explains the weak
inequality in Corollary~\ref{u:cor:closure}, without changing any
finite-stage determinant identity.

\subsection{What the rank laws determine}
The exact count $M_\tau(\varepsilon)$ is an existence and obstruction
law. It does not bound the size or conductance ratios of a graph
approximating a prescribed extremal spectrum. The degree-three
construction gives finite examples but can use very large reservoirs
and widely separated sizes. The centered-path coefficient $1/(2\pi^2)$ also equals the uniform
worst-case leading coefficient over all finite unit paths and cycles.
The matching upper uses a fixed continuum operator dominating every
finite low-frequency grid, rather than inferring uniformity from
fixed-index spectral convergence.

The constructive theorem gives explicit two-sided relative factors
from graph primitives. Its near-sharp rank variant can approach the
structural leading rank constant, but it does not compute the exact
one-sided spectral truncation. A generalized Nystr\"om alternative
can also be analyzed using ridge effective dimension and classical
Gaussian estimates~\cite{FTU}; its fractional sampling and matrix-action
oracles must be fully charged. The flow-based construction avoids
treating those fractional vector operations as free.

\subsection{Relation to matrix-function approximation}
The rational rank identity used in an alternative degree-two order bound
belongs to classical rational matrix-function updating; Beckermann, Cortinovis,
Kressner and Schweitzer discuss its connection to Bernstein and
Van Loan's rank-one identities~\cite[Section 3.3]{BCKS}.
Their Theorem 4.5 bounds an ordinary operator-norm update error using
uniform scalar rational approximation. Converting such an error $E$
to the present metric by
\[
 \norm{(L^\dagger)^{\alpha/2}E(L^\dagger)^{\alpha/2}}
 \le\lambda_2(L)^{-\alpha}\norm E
\]
can reintroduce a graph spectral gap. The square-root Riccati framework
of Shmueli, Drineas and Avron supplies update decay and error analysis,
with costs that explicitly charge the supplied square-root and inverse
square-root actions~\cite[Section 5.5]{SDA}. Those are established
update methods with different error or access assumptions.

Contour and sinc quadrature for matrix fractional powers precede the
present construction~\cite{HHT,BLP,TW}. Bonito, Lei and Pasciak also
analyze fractional operator-domain norms~\cite[Theorem 3.2]{BLP}.
Benzi and Simunec study rational Krylov computation of fractional
diffusion vectors on graphs~\cite{BS}. Persson, Meyer and Musco transfer
low-rank approximations to an operator-monotone function of a single
matrix~\cite{PMM}. Applying a low-rank method directly to the present
$C$ still requires a charged representation or action for that matrix.
The additional graph-specific steps here are the dipole moment bounds,
feasible-flow conversion to fractional relative accuracy, and the
computable relative Gram. These comparisons distinguish statements;
they are not an exhaustive historical priority assessment.

Perturbation determinants and their trace formulas have extensive
classical theory~\cite{GZ,Bornemann}. Determinants, local graph limits,
and resistance also meet in Lyons's work on tree entropy~\cite{LyonsEntropy}.
That normalized one-operator entropy differs from the unnormalized
relative fractional determinant for a distinguished edge used here.
Kron reduction and resistance preservation are established circuit
principles~\cite{DorflerBullo}, and regular-tree capacities and free/wired
resistance conventions are classical network theory~\cite{LyonsPeres}.
The sharp graph count combines these inputs with a multiport realization;
none of these underlying mechanisms is claimed as a new general theory.

The fixed-degree outlier and the arithmetic theorem address separate
constraints. The former limits possible spectra. The latter bounds
the cost of approximating a given graph response in a stated arithmetic
model. Sharp leading constants are proved at every fixed degree $d\ge2$.
No effective rate is supplied for the compactness-based higher-degree
asymptotic. Finite-word stability and historical priority remain
unresolved here.
